% Part: normal-modal-logic
% Chapter: axioms-systems
% Section: logics-proofs

\documentclass[../../../include/open-logic-section]{subfiles}

\begin{document}

\olfileid[pt]{nml}{prf}{prf}

\olsection{\usetoken{P}{derivation} e sistemas modais}

Primeiro definimos o que é !!a{derivation} para lógicas modais
normais. Em linhas gerais, !!a{derivation} é uma sequência de
!!{formula}s na qual cada !!{element} é ou (uma instância de
substituição de) um de certos \emph{axiomas}, ou decorre de
!!{element}s anteriores por uma de algumas regras de inferência. Para
lógicas modais normais, todas as instâncias de
tautologias\iftag{prvDiamond}{, \Ax{K} e~\Dual{}}{ e~\Ax{K}}
contam como axiomas. Isso resulta no sistema modal~$\Log{K}$, a menor
lógica modal normal. Podemos, contudo, querer acrescentar axiomas para
obter outros sistemas. As regras são sempre modus ponens~\MP{} e
necessitação~\Nec.

\begin{defn}
  Dado um sistema modal $\Log{K} !A_1 \dots !A_n$ e !!a{formula}~$!B$,
  dizemos que $!B$ é \emph{!!{derivable}} em $\Log{K} !A_1 \dots
  !A_n$, e escrevemos $\Log{K} !A_1 \dots !A_n \Proves !B$, se e
  somente se existem !!{formula}s $!C_1$, \dots, $!C_k$ tais que
  $!C_k = !B$ e cada $!C_i$ é ou uma instância tautológica, ou uma
  instância de um dos axiomas $\Ax{K}$,\iftag{prvDiamond}{ $\Dual$,}{}
  $!A_1$, \dots, $!A_n$, ou decorre de !!{formula}s anteriores por
  meio das regras \MP{} ou~\Nec.
\end{defn}

A proposição seguinte nos permite mostrar que $!B \in \Sigma$
exibindo uma $\Sigma$-!!{derivation} de~$!B$.

\begin{prop}
  $\Log{K} !A_1 \dots !A_n = \Setabs{!B}{\Log{K} !A_1 \dots !A_n \Proves !B}$.
\end{prop}

\begin{proof}
  Usamos indução no comprimento de !!{derivation}s para mostrar que
  $\Setabs{!B}{\Log{K} !A_1 \dots !A_n \Proves !B} \subseteq
  \Log{K} !A_1 \dots !A_n$.

  Se a !!{derivation} de~$!B$ tem comprimento~$1$, ela contém uma única
  !!{formula}. Essa !!{formula} não pode decorrer de fórmulas anteriores
  por \MP{} ou \Nec, portanto deve ser uma instância tautológica, uma
  instância de \Ax{K},\iftag{prvDiamond}{ $\Dual$,}{} ou uma instância
  de um dos axiomas $!A_1$, \dots,~$!A_n$. Mas
  $\Log{K}!A_1\dots!A_n$ também contém essas instâncias, logo
  $!B \in \Log{K}!A_1 \dots !A_n$.

  Se a !!{derivation} de $!B$ tem comprimento~$> 1$, então $!B$ também
  pode ser obtida por \MP{} ou \Nec{} a partir de !!{formula}s que não
  ocorrem como a última linha da !!{derivation}. Se $!B$ decorre de
  $!C$ e $!C \lif !B$ (por \MP), então $!C$ e $!C \lif !B \in
  \Log{K}!A_1 \dots !A_n$ pela hipótese de indução. Mas toda lógica
  modal é fechada sob modus ponens, logo $!B \in \Log{K}!A_1 \dots
  !A_n$. Se $!B \equiv \Box !C$ decorre de $!C$ por \Nec, então $!C
  \in \Log{K}!A_1 \dots !A_n$ pela hipótese de indução. Mas toda
  lógica modal normal é fechada sob $\Nec$, logo $!B \in
  \Log{K}!A_1\dots!A_n$.

  A inclusão inversa resulta de mostrar que
  $\Sigma = \Setabs{!B}{\Log{K} !A_1 \dots !A_n \Proves !B }$ é uma
  lógica modal normal que contém todas as instâncias de $!A_1$,
  \dots, $!A_n$, e de observar que $\Log{K} !A_1 \dots !A_n$ é, por
  definição, a menor lógica desse tipo.
  \begin{enumerate}
    \item Toda tautologia~$!B$ é uma instância tautológica, portanto
      $\Log{K}!A_1\dots!A_n \Proves !B$, e assim $\Sigma$ contém todas
      as tautologias.
    \item Se $\Log{K}!A_1\dots!A_n \Proves !C$ e
      $\Log{K}!A_1\dots!A_n \Proves !C \lif !B$, então
      $\Log{K}!A_1\dots!A_n \Proves !B$: combinamos a !!{derivation} de
      $!C$ com a de $!C \lif !B$ e acrescentamos a linha~$!B$. A última
      linha é justificada por \MP{}. Portanto, $\Sigma$ é fechada sob
      modus ponens.
    \item Se $!B$ tem !!a{derivation}, então toda instância de
      substituição de~$!B$ também tem !!a{derivation}: aplicamos a
      substituição a cada !!{formula} da !!{derivation}. (Exercício:
      prove, por indução no comprimento de !!{derivation}s, que o
      resultado também é uma !!{derivation} correta.) Portanto,
      $\Sigma$ é fechada sob substituição uniforme. (Já estabelecemos
      que $\Sigma$ satisfaz todas as condições de uma lógica modal.)
    \item Temos $\Log{K}!A_1\dots!A_n \Proves \Ax{K}$, logo $K \in \Sigma$.
    \tagitem{prvDiamond}{Temos
      $\Log{K}!A_1\dots!A_n \Proves \Dual$, logo $\Dual \in \Sigma$.}{}
    \item Se $\Log{K}!A_1\dots!A_n \Proves !C$, a linha adicional
      $\Box !C$ é justificada por~\Nec. Consequentemente, $\Sigma$ é
      fechada sob~\Nec. Assim, $\Sigma$ é normal.
  \end{enumerate}
\end{proof}

\end{document}
